
\documentclass{article}
\usepackage{colortbl}
\usepackage{makecell}
\usepackage{multirow}
\usepackage{supertabular}

\begin{document}

\newcounter{utterance}

\centering \large Interaction Transcript for game `hot\_air\_balloon', experiment `air\_balloon\_survival\_en\_complexity\_easy', episode 5 with deepseek{-}chat{-}v3.1{-}t1.0.
\vspace{24pt}

{ \footnotesize  \setcounter{utterance}{1}
\setlength{\tabcolsep}{0pt}
\begin{supertabular}{c@{$\;$}|p{.15\linewidth}@{}p{.15\linewidth}p{.15\linewidth}p{.15\linewidth}p{.15\linewidth}p{.15\linewidth}}
    \# & $\;$A & \multicolumn{4}{c}{Game Master} & $\;\:$B\\
    \hline

    \theutterance \stepcounter{utterance}  
    & & \multicolumn{4}{p{0.6\linewidth}}{
        \cellcolor[rgb]{0.9,0.9,0.9}{
            \makecell[{{p{\linewidth}}}]{
                \texttt{\tiny{[P1$\langle$GM]}}
                \texttt{Stai partecipando a un gioco di negoziazione collaborativa.} \\
\\ 
\texttt{Insieme a un altro partecipante, dovete accordarvi su un unico set di oggetti da conservare. Ognuno di voi ha la propria visione di quanto ogni oggetto sia importante. Non sai come l'altro partecipante valuta gli oggetti. Inoltre, ti viene fornito lo sforzo che ogni oggetto richiede.} \\
\texttt{Potete accordarvi su un set solo se lo sforzo totale degli oggetti selezionati non supera un limite condiviso:} \\
\\ 
\texttt{LIMITE = 9475} \\
\\ 
\texttt{Ecco i valori di sforzo dei singoli oggetti:} \\
\\ 
\texttt{Sforzo oggetto = \{"B21": 944, "B00": 568, "A08": 515, "A27": 239, "B94": 876, "B53": 421, "A23": 991, "C65": 281, "A89": 788, "B67": 678, "B45": 432, "C34": 409, "C12": 279, "A05": 506, "B60": 101, "C92": 686, "B69": 851, "A84": 854, "C06": 133, "B12": 192, "A16": 573, "B85": 17, "C35": 465, "B10": 772, "C15": 46, "B09": 501, "B16": 220, "C81": 404, "A58": 843, "C24": 748, "A30": 552, "A66": 843, "B75": 947, "B46": 345, "C70": 930\}} \\
\\ 
\texttt{Ecco la tua visione personale sull'importanza di ogni oggetto:} \\
\\ 
\texttt{Valori di importanza dell'oggetto = \{"B21": 138, "B00": 583, "A08": 868, "A27": 822, "B94": 783, "B53": 65, "A23": 262, "C65": 121, "A89": 508, "B67": 780, "B45": 461, "C34": 484, "C12": 668, "A05": 389, "B60": 808, "C92": 215, "B69": 97, "A84": 500, "C06": 30, "B12": 915, "A16": 856, "B85": 400, "C35": 444, "B10": 623, "C15": 781, "B09": 786, "B16": 3, "C81": 713, "A58": 457, "C24": 273, "A30": 739, "A66": 822, "B75": 235, "B46": 606, "C70": 968\}} \\
\\ 
\texttt{Obiettivo:} \\
\\ 
\texttt{Il tuo obiettivo è negoziare un set condiviso di oggetti che ti avvantaggi il più possibile (cioè, massimizzi l'importanza totale per TE), rimanendo entro il LIMITEE. Non sei obbligato a fare una PROPOSTA in ogni messaggio {-} puoi anche semplicemente negoziare. Tutte le tattiche sono permesse!} \\
\\ 
\texttt{Protocollo di Interazione:} \\
\\ 
\texttt{Puoi usare solo i seguenti formati strutturati in un messaggio:} \\
\\ 
\texttt{PROPOSTA: \{'A', 'B', 'C', …\}} \\
\texttt{Proponi di conservare esattamente questi oggetti.} \\
\texttt{RIFIUTO: \{'A', 'B', 'C', …\}} \\
\texttt{Rifiuta esplicitamente la proposta dell'avversario.} \\
\texttt{ARGOMENTAZIONE: \{'...'\}} \\
\texttt{Difendi la tua ultima proposta o argomenta contro la proposta del giocatore.} \\
\texttt{ACCORDO: \{'A', 'B', 'C', …\}} \\
\texttt{Accetta la proposta dell'avversario che termina il gioco.} \\
\texttt{RAGIONAMENTO STRATEGICO: \{'...'\}} \\
\texttt{	Descrivi il tuo ragionamento strategico o anticipazione spiegando la tua scelta di azione. Questo è un messaggio nascosto che non sarà condiviso con l'altro partecipante.} \\
\\ 
\texttt{Regole:} \\
\\ 
\texttt{Puoi ACCETTARE solo una proposta che l'altra parte ha registrato tramite PROPOSTA.} \\
\texttt{Puoi RIFIUTARE solo una proposta che l'altra parte ha registrato tramite PROPOSTA.} \\
\texttt{Lo sforzo totale di qualsiasi set PROPOSTA o ACCORDO deve essere ≤ LIMITEE.} \\
\texttt{NON rivelare i tuoi punteggi di importanza nascosti.} \\
\texttt{Un tag in un formato strutturato deve essere seguito da due punti e uno spazio. L'argomento deve essere un set python contenente 0 o più stringhe.} \\
\texttt{Quindi, deve essere della forma TAG: \{...\}} \\
\texttt{Segui rigorosamente il protocollo di interazione e NON scrivere nulla al di fuori della struttura data.} \\
\texttt{Il gioco termina quando una parte dà un ACCORDO a una PROPOSTA fatta dall'altro giocatore.} \\
\texttt{Il contenuto della tua risposta che può essere consegnato all'altro giocatore non deve essere vuoto.} \\
\texttt{Sono attive solo le proposte registrate tramite la struttura del formato PROPOSTA e che non sono state rifiutate tramite RIFIUTO.} \\
\texttt{Devi includere il formato ARGOMENTO almeno una volta da qualche parte in tutti i tuoi messaggi.} \\
\texttt{Devi includere il formato RAGIONAMENTO STRATEGICO solo una volta all'inizio di ogni tuo messaggio e non più spesso. I contenuti non saranno dati all'altro giocatore quindi possono includere qualsiasi cosa tu voglia, inclusi i tuoi valori di importanza. Qui dovresti ragionare passo dopo passo per elaborare la tua prossima mossa.} \\
            }
        }
    }
    & \\ \\

    \theutterance \stepcounter{utterance}  
    & \multicolumn{4}{p{0.6\linewidth}}{
        \cellcolor[rgb]{0.9,0.9,0.9}{
            \makecell[{{p{\linewidth}}}]{
                \texttt{\tiny{[P1$\rangle$GM]}}
                \texttt{RAGIONAMENTO STRATEGICO: Devo massimizzare il mio valore totale (138+583+868+822+783+65+262+121+508+780+461+484+668+389+808+215+97+500+30+915+856+400+444+623+781+786+3+713+457+273+739+822+235+606+968 = 20213) rimanendo sotto il limite di sforzo 9475. Analizzo il rapporto valore/sforzo per prioritizzare gli oggetti. I migliori sono: B85 (400/17=23.53), C15 (781/46=16.98), B60 (808/101=8.0), B12 (915/192=4.77), A08 (868/515=1.69), A27 (822/239=3.44), B67 (780/678=1.15), C12 (668/279=2.39), A16 (856/573=1.49), B10 (623/772=0.81). Considerando che l'altro giocatore ha valori diversi, devo negoziare un set che massimizzi il mio valore mentre convince l'altro. Inizierò proponendo gli oggetti con alto rapporto valore/sforzo per me.} \\
\\ 
\texttt{PROPOSTA: \{'B85', 'C15', 'B60', 'B12', 'A08', 'A27', 'B67', 'C12', 'A16', 'B10'\}} \\
\texttt{ARGOMENTAZIONE: Questa proposta include oggetti a basso sforzo e alto valore, efficienti per entrambi. Lo sforzo totale è 17+46+101+192+515+239+678+279+573+772 = 3412, ben sotto il limite, lasciando spazio per aggiunte collaborative.} \\
            }
        }
    }
    & & \\ \\

    \theutterance \stepcounter{utterance}  
    & & \multicolumn{4}{p{0.6\linewidth}}{
        \cellcolor[rgb]{0.9,0.9,0.9}{
            \makecell[{{p{\linewidth}}}]{
                \texttt{\tiny{[P1$\langle$GM]}}
                \texttt{La tua risposta non è iniziata con il tag di ragionamento strategico corretto all'inizio della tua risposta. Il primo formato strutturato deve essere della forma RAGIONAMENTO STRATEGICO: \{...\}. Riprova.} \\
            }
        }
    }
    & \\ \\

    \theutterance \stepcounter{utterance}  
    & \multicolumn{4}{p{0.6\linewidth}}{
        \cellcolor[rgb]{0.9,0.9,0.9}{
            \makecell[{{p{\linewidth}}}]{
                \texttt{\tiny{[P1$\rangle$GM]}}
                \texttt{RAGIONAMENTO STRATEGICO: \{'Devo massimizzare il mio valore totale (calcolato approssimativamente 20213) rimanendo sotto il limite di sforzo 9475. Analizzo il rapporto valore/sforzo per prioritizzare gli oggetti. I migliori sono: B85 (400/17=23.53), C15 (781/46=16.98), B60 (808/101=8.0), B12 (915/192=4.77), A27 (822/239=3.44), C12 (668/279=2.39), A08 (868/515=1.69), A16 (856/573=1.49), B67 (780/678=1.15), B10 (623/772=0.81). Considerando che l\textbackslash{}'altro giocatore ha valori diversi, devo negoziare un set che massimizzi il mio valore mentre è attraente per l\textbackslash{}'altro. Inizierò proponendo un set con oggetti ad alto rapporto per me, ma con spazio per aggiunte.'\}} \\
\\ 
\texttt{PROPOSTA: \{'B85', 'C15', 'B60', 'B12', 'A27', 'C12', 'A08', 'A16', 'B67', 'B10'\}} \\
\texttt{ARGOMENTAZIONE: \{'Questa proposta include oggetti a basso sforzo e alto valore, molto efficienti per entrambi. Lo sforzo totale è 17+46+101+192+239+279+515+573+678+772 = 3412, ben sotto il limite di 9475, lasciando ampio spazio per aggiungere altri oggetti che potrebbero essere importanti per te.'\}} \\
            }
        }
    }
    & & \\ \\

    \theutterance \stepcounter{utterance}  
    & & \multicolumn{4}{p{0.6\linewidth}}{
        \cellcolor[rgb]{0.9,0.9,0.9}{
            \makecell[{{p{\linewidth}}}]{
                \texttt{\tiny{[GM$\rangle$P2]}}
                \texttt{Stai partecipando a un gioco di negoziazione collaborativa.} \\
\\ 
\texttt{Insieme a un altro partecipante, dovete accordarvi su un unico set di oggetti da conservare. Ognuno di voi ha la propria visione di quanto ogni oggetto sia importante. Non sai come l'altro partecipante valuta gli oggetti. Inoltre, ti viene fornito lo sforzo che ogni oggetto richiede.} \\
\texttt{Potete accordarvi su un set solo se lo sforzo totale degli oggetti selezionati non supera un limite condiviso:} \\
\\ 
\texttt{LIMITE = 9475} \\
\\ 
\texttt{Ecco i valori di sforzo dei singoli oggetti:} \\
\\ 
\texttt{Sforzo oggetto = \{"B21": 944, "B00": 568, "A08": 515, "A27": 239, "B94": 876, "B53": 421, "A23": 991, "C65": 281, "A89": 788, "B67": 678, "B45": 432, "C34": 409, "C12": 279, "A05": 506, "B60": 101, "C92": 686, "B69": 851, "A84": 854, "C06": 133, "B12": 192, "A16": 573, "B85": 17, "C35": 465, "B10": 772, "C15": 46, "B09": 501, "B16": 220, "C81": 404, "A58": 843, "C24": 748, "A30": 552, "A66": 843, "B75": 947, "B46": 345, "C70": 930\}} \\
\\ 
\texttt{Ecco la tua visione personale sull'importanza di ogni oggetto:} \\
\\ 
\texttt{Valori di importanza dell'oggetto = \{"B21": 138, "B00": 583, "A08": 868, "A27": 822, "B94": 783, "B53": 65, "A23": 262, "C65": 121, "A89": 508, "B67": 780, "B45": 461, "C34": 484, "C12": 668, "A05": 389, "B60": 808, "C92": 215, "B69": 97, "A84": 500, "C06": 30, "B12": 915, "A16": 856, "B85": 400, "C35": 444, "B10": 623, "C15": 781, "B09": 786, "B16": 3, "C81": 713, "A58": 457, "C24": 273, "A30": 739, "A66": 822, "B75": 235, "B46": 606, "C70": 968\}} \\
\\ 
\texttt{Obiettivo:} \\
\\ 
\texttt{Il tuo obiettivo è negoziare un set condiviso di oggetti che ti avvantaggi il più possibile (cioè, massimizzi l'importanza totale per TE), rimanendo entro il LIMITEE. Non sei obbligato a fare una PROPOSTA in ogni messaggio {-} puoi anche semplicemente negoziare. Tutte le tattiche sono permesse!} \\
\\ 
\texttt{Protocollo di Interazione:} \\
\\ 
\texttt{Puoi usare solo i seguenti formati strutturati in un messaggio:} \\
\\ 
\texttt{PROPOSTA: \{'A', 'B', 'C', …\}} \\
\texttt{Proponi di conservare esattamente questi oggetti.} \\
\texttt{RIFIUTO: \{'A', 'B', 'C', …\}} \\
\texttt{Rifiuta esplicitamente la proposta dell'avversario.} \\
\texttt{ARGOMENTAZIONE: \{'...'\}} \\
\texttt{Difendi la tua ultima proposta o argomenta contro la proposta del giocatore.} \\
\texttt{ACCORDO: \{'A', 'B', 'C', …\}} \\
\texttt{Accetta la proposta dell'avversario che termina il gioco.} \\
\texttt{RAGIONAMENTO STRATEGICO: \{'...'\}} \\
\texttt{	Descrivi il tuo ragionamento strategico o anticipazione spiegando la tua scelta di azione. Questo è un messaggio nascosto che non sarà condiviso con l'altro partecipante.} \\
\\ 
\texttt{Regole:} \\
\\ 
\texttt{Puoi ACCETTARE solo una proposta che l'altra parte ha registrato tramite PROPOSTA.} \\
\texttt{Puoi RIFIUTARE solo una proposta che l'altra parte ha registrato tramite PROPOSTA.} \\
\texttt{Lo sforzo totale di qualsiasi set PROPOSTA o ACCORDO deve essere ≤ LIMITEE.} \\
\texttt{NON rivelare i tuoi punteggi di importanza nascosti.} \\
\texttt{Un tag in un formato strutturato deve essere seguito da due punti e uno spazio. L'argomento deve essere un set python contenente 0 o più stringhe.} \\
\texttt{Quindi, deve essere della forma TAG: \{...\}} \\
\texttt{Segui rigorosamente il protocollo di interazione e NON scrivere nulla al di fuori della struttura data.} \\
\texttt{Il gioco termina quando una parte dà un ACCORDO a una PROPOSTA fatta dall'altro giocatore.} \\
\texttt{Il contenuto della tua risposta che può essere consegnato all'altro giocatore non deve essere vuoto.} \\
\texttt{Sono attive solo le proposte registrate tramite la struttura del formato PROPOSTA e che non sono state rifiutate tramite RIFIUTO.} \\
\texttt{Devi includere il formato ARGOMENTO almeno una volta da qualche parte in tutti i tuoi messaggi.} \\
\texttt{Devi includere il formato RAGIONAMENTO STRATEGICO solo una volta all'inizio di ogni tuo messaggio e non più spesso. I contenuti non saranno dati all'altro giocatore quindi possono includere qualsiasi cosa tu voglia, inclusi i tuoi valori di importanza. Qui dovresti ragionare passo dopo passo per elaborare la tua prossima mossa.} \\
\\ 
\texttt{Ora riceverai il primo messaggio dell'altro giocatore.} \\
\\ 
\texttt{PROPOSTA: \{'B85', 'C15', 'B60', 'B12', 'A27', 'C12', 'A08', 'A16', 'B67', 'B10'\}} \\
\texttt{ARGOMENTAZIONE: \{'Questa proposta include oggetti a basso sforzo e alto valore, molto efficienti per entrambi. Lo sforzo totale è 17+46+101+192+239+279+515+573+678+772 = 3412, ben sotto il limite di 9475, lasciando ampio spazio per aggiungere altri oggetti che potrebbero essere importanti per te.'\}} \\
            }
        }
    }
    & \\ \\

    \theutterance \stepcounter{utterance}  
    & & & \multicolumn{4}{p{0.6\linewidth}}{
        \cellcolor[rgb]{0.9,0.9,0.9}{
            \makecell[{{p{\linewidth}}}]{
                \texttt{\tiny{[GM$\langle$P2]}}
                \texttt{RAGIONAMENTO STRATEGICO: \{"Analizzo la proposta iniziale dell'avversario. Lo sforzo totale è 3412, ben sotto il limite. I miei oggetti proposti: B85 (valore 400), C15 (781), B60 (808), B12 (915), A27 (822), C12 (668), A08 (868), A16 (856), B67 (780), B10 (623). Il valore totale per me sarebbe 400+781+808+915+822+668+868+856+780+623 = 7521. Questa è un'ottima base, ma posso migliorarla aggiungendo altri oggetti ad alto valore per me. I miei oggetti migliori non inclusi sono: C70 (968), B00 (583), A30 (739), A66 (822), B45 (461), C34 (484), C81 (713), A89 (508), A58 (457), B46 (606). Dovrei proporre un set più completo che massimizzi il mio valore rimanendo sotto il limite. Lo sforzo della proposta corrente è basso, quindi posso aggiungere molti oggetti. Calcolo uno set ideale per me: parto dalla loro proposta e aggiungo i miei oggetti migliori fino ad avvicinarmi al limite. Sforzo corrente: 3412. Posso aggiungere: C70 (sforzo 930, totale 4342), A30 (552, totale 4894), A66 (843, totale 5737), B45 (432, totale 6169), C34 (409, totale 6578), C81 (404, totale 6982), A89 (788, totale 7770), A58 (843, totale 8613), B46 (345, totale 8958), e forse altri. Totale sforzo 8958, valore aggiunto 968+739+822+461+484+713+508+457+606=5758, valore totale 7521+5758=13279. Questo è eccellente. Propongo questo set ampliato."\}} \\
\\ 
\texttt{PROPOSTA: \{'B85', 'C15', 'B60', 'B12', 'A27', 'C12', 'A08', 'A16', 'B67', 'B10', 'C70', 'A30', 'A66', 'B45', 'C34', 'C81', 'A89', 'A58', 'B46'\}} \\
\texttt{ARGOMENTAZIONE: \{'Accetto la tua proposta di base perché è efficiente, ma suggerisco di espanderla per includere più oggetti ad alto valore per entrambi, dato l'ampio margine disponibile. Lo sforzo totale di questo set è 8958, ben sotto il limite di 9475, quindi è fattibile. Questo set cattura molti oggetti con buon rapporto valore/sforzo.'\}} \\
            }
        }
    }
     \\ \\

    \theutterance \stepcounter{utterance}  
    & & \multicolumn{4}{p{0.6\linewidth}}{
        \cellcolor[rgb]{0.9,0.9,0.9}{
            \makecell[{{p{\linewidth}}}]{
                \texttt{\tiny{[P1$\langle$GM]}}
                \texttt{PROPOSTA: \{'B85', 'C15', 'B60', 'B12', 'A27', 'C12', 'A08', 'A16', 'B67', 'B10', 'C70', 'A30', 'A66', 'B45', 'C34', 'C81', 'A89', 'A58', 'B46'\}} \\
\texttt{ARGOMENTAZIONE: \{'Accetto la tua proposta di base perché è efficiente, ma suggerisco di espanderla per includere più oggetti ad alto valore per entrambi, dato l'ampio margine disponibile. Lo sforzo totale di questo set è 8958, ben sotto il limite di 9475, quindi è fattibile. Questo set cattura molti oggetti con buon rapporto valore/sforzo.'\}} \\
            }
        }
    }
    & \\ \\

    \theutterance \stepcounter{utterance}  
    & \multicolumn{4}{p{0.6\linewidth}}{
        \cellcolor[rgb]{0.9,0.9,0.9}{
            \makecell[{{p{\linewidth}}}]{
                \texttt{\tiny{[P1$\rangle$GM]}}
                \texttt{RAGIONAMENTO STRATEGICO: \{'La proposta dell\textbackslash{}'altro giocatore aggiunge molti oggetti: C70 (sforzo 930, mio valore 968), A30 (552, 739), A66 (843, 822), B45 (432, 461), C34 (409, 484), C81 (404, 713), A89 (788, 508), A58 (843, 457), B46 (345, 606). Il mio valore totale per questo set sarebbe: 400+781+808+915+822+668+868+856+780+623+968+739+822+461+484+713+508+457+606 = 13879. Lo sforzo totale è 3412 + (930+552+843+432+409+404+788+843+345) = 3412 + 5546 = 8958, sotto il limite. Questo è un ottimo valore per me, ma posso cercare di aggiungere altri oggetti ad alto rapporto come B09 (786/501=1.57) o C35 (444/465=0.95) per massimizzare ulteriormente, dato che c\textbackslash{}'è ancora margine (9475{-}8958=517). Tuttavia, poiché l\textbackslash{}'altro sembra voler collaborare, accetterei questa proposta ricca.'\}} \\
\\ 
\texttt{ACCORDO: \{'B85', 'C15', 'B60', 'B12', 'A27', 'C12', 'A08', 'A16', 'B67', 'B10', 'C70', 'A30', 'A66', 'B45', 'C34', 'C81', 'A89', 'A58', 'B46'\}} \\
            }
        }
    }
    & & \\ \\

    \theutterance \stepcounter{utterance}  
    & & \multicolumn{4}{p{0.6\linewidth}}{
        \cellcolor[rgb]{0.9,0.9,0.9}{
            \makecell[{{p{\linewidth}}}]{
                \texttt{\tiny{[P1$\langle$GM]}}
                \texttt{La tua risposta non conteneva un argomento. Devi includere il formato strutturato ARGOMENTO: \{...\} da qualche parte nella tua risposta. Riprova.} \\
            }
        }
    }
    & \\ \\

    \theutterance \stepcounter{utterance}  
    & \multicolumn{4}{p{0.6\linewidth}}{
        \cellcolor[rgb]{0.9,0.9,0.9}{
            \makecell[{{p{\linewidth}}}]{
                \texttt{\tiny{[P1$\rangle$GM]}}
                \texttt{RAGIONAMENTO STRATEGICO: \{'La proposta dell\textbackslash{}'altro giocatore è molto vantaggiosa per me, con un valore totale stimato di oltre 13800 e uno sforzo di 8958, ben sotto il limite. Include quasi tutti i miei oggetti ad alto valore e alto rapporto. L\textbackslash{}'unico oggetto significativo che manca è B09 (valore 786, sforzo 501), ma lo sforzo residuo di 517 sarebbe sufficiente per aggiungerlo. Tuttavia, data l\textbackslash{}'ampiezza della proposta e il suo alto valore, accettarla direttamente è la scelta più strategica per concludere rapidamente un accordo vantaggioso senza rischiare un rifiuto.'\}} \\
\\ 
\texttt{ARGOMENTAZIONE: \{'Questa è una proposta eccellente e collaborativa. Cattura una vasta gamma di oggetti ad alto valore per entrambi, utilizzando in modo efficiente la maggior parte del budget disponibile. Sono d\textbackslash{}'accordo sul fatto che sia un set bilanciato e vantaggioso.'\}} \\
\\ 
\texttt{ACCORDO: \{'B85', 'C15', 'B60', 'B12', 'A27', 'C12', 'A08', 'A16', 'B67', 'B10', 'C70', 'A30', 'A66', 'B45', 'C34', 'C81', 'A89', 'A58', 'B46'\}} \\
            }
        }
    }
    & & \\ \\

    \theutterance \stepcounter{utterance}  
    & & & \multicolumn{2}{p{0.3\linewidth}}{
        \cellcolor[rgb]{0.9,0.9,0.9}{
            \makecell[{{p{\linewidth}}}]{
                \texttt{\tiny{[GM$|$GM]}}
                \texttt{game successful} \\
            }
        }
    }
    & & \\ \\

    \theutterance \stepcounter{utterance}  
    & & & \multicolumn{2}{p{0.3\linewidth}}{
        \cellcolor[rgb]{0.9,0.9,0.9}{
            \makecell[{{p{\linewidth}}}]{
                \texttt{\tiny{[GM$|$GM]}}
                \texttt{end game} \\
            }
        }
    }
    & & \\ \\

\end{supertabular}
}

\end{document}
